\documentclass[11pt]{article}
\usepackage[a4paper,margin=27mm]{geometry}
\usepackage{amsmath,amssymb,amsthm,hyperref,booktabs}
\newtheorem{proposition}{Proposition}
\newtheorem{lemma}{Lemma}
\newcommand{\FPdim}{\operatorname{FPdim}}
\title{A small relative-center target for the last multiplicity-eight ring}
\author{}\date{30 September 2026}
\begin{document}\maketitle
The ring $R6_{8,1660}$ remains unresolved. The following reduction uses its
Ising subcategory and module-category rigidity. It does not enumerate induction
matrices or modular data for the full Drinfeld center.

Write its basis as $1,g,X,gX,\sigma,\sigma X$. Its multiplication is determined by
\[
 g^2=1,\qquad g\sigma=\sigma,\qquad \sigma^2=1+g,\qquad
 X^2=1+\bigl(3(1+g)+4\sigma\bigr)X,
\]
with all simples self-dual and the displayed generators commuting.
Set $d=3+2\sqrt2+2\sqrt3+\sqrt6$. Then
$\FPdim X=d$ and $\FPdim\mathcal C=4(1+d^2)$.

\begin{lemma}
For a fusion subcategory $\mathcal D\subset\mathcal C$, the category
$Z(\mathcal D)$ has a central lift, with reverse braiding, to
$Z(Z_{\mathcal D}(\mathcal C))$. In particular its Grothendieck ring
is central in that of the relative center.
\end{lemma}
\begin{proof}
Use strict notation and write a relative-central object as $(X,\gamma)$,
where $\gamma_V:VX\to XV$ for $V\in\mathcal D$. For
$b=(Y,\beta)\in Z(\mathcal D)$, define its commutation with $(X,\gamma)$
by $\gamma_Y:YX\to XY$. To check that this is a relative-central morphism,
we must verify
\[
 (\gamma_Y\otimes1_V)(1_Y\otimes\gamma_V)(\beta_V\otimes1_X)
 =(1_X\otimes\beta_V)(\gamma_V\otimes1_Y)(1_V\otimes\gamma_Y).
\]
The two half-braiding identities rewrite this as naturality of $\gamma$
at the morphism $\beta_V:VY\to YV$. Naturality in relative-central objects
and tensor compatibility follow from their defining identities. On
$Z(\mathcal D)$ the resulting braiding is the reverse of its usual braiding.
\end{proof}

\begin{proposition}
Any categorification $\mathcal C$ of $R6_{8,1660}$ has a rank-eighteen
relative center $\mathcal E=Z_{\mathcal I}(\mathcal C)$, where
$\mathcal I=\langle g,\sigma\rangle$ is Ising. Write
$B=K_0(Z(\mathcal I))\cong K_0(\mathcal I)\otimes K_0(\mathcal I)$ and
$e=g\otimes g$. Its Grothendieck ring has basis $b,bt$ for the nine
basis elements $b$ of $B$, with
\[
 tb=bt,\qquad t^2=L+At,\qquad L\in\{1,e\},\qquad A\in B_{\ge0},
 \qquad A=LA.
\]
The forgetful map on $B$ is $F(a\otimes b)=ab$, and
\[
 F(t)=X,\qquad F(A)=3(1+g)+4\sigma,\qquad t^*=Lt.
\]
\end{proposition}
\begin{proof}
The complement $\mathcal M=\langle X,gX,\sigma X\rangle$ is an
$\mathcal I$-bimodule category. The functor $Y\mapsto YX$ is a left
module equivalence from $\mathcal I$ onto $\mathcal M$: it takes the
three simples bijectively to the three simples, or equivalently
$\underline{\operatorname{End}}_{\mathcal I}(X)=1$.
Transport the right action to this regular left module. It becomes a
tensor autoequivalence of $\mathcal I^{\rm rev}$, since it fixes all
simple classes. Every tensor autoequivalence of an Ising category is
isomorphic to the identity \cite[Remark B.6(i)]{DGNO}. Thus
$\mathcal M$ is the regular bimodule, with $1$ corresponding to $X$.

Passing to bimodule centers gives a regular $Z(\mathcal I)$-module
$Z_{\mathcal I}(\mathcal M)$ \cite[Definition 2.1 and Remark 2.2]{GNN}.
Consequently its simples are $bt$, and $\mathcal E$ has eighteen
simples. The lemma gives commutativity. The part of $t^2$ supported on
$Z(\mathcal I)$ forgets to $1$, and hence is exactly $1$ or $e$.
Frobenius reciprocity gives $t^*=Lt$. All simples of $B$ are self-dual;
taking duals of $t^2=L+At$ therefore gives $A=LA$. The remaining
statements follow by applying the forgetful functor.
\end{proof}

\begin{proposition}
Using the previously proved reduction to a pseudo-unitary
categorification, the above relative-center rings reduce to $58$
labelled alternatives, covered by fourteen explicitly listed based-ring
representatives.
\end{proposition}
\begin{proof}
The relative center inherits positive spherical dimensions from
$\mathcal C$. Its dimension and dimension field are
\[
 D=16(1+d^2)=576+384\sqrt2+320\sqrt3+224\sqrt6,
 \qquad K=\mathbb Q(\sqrt2,\sqrt3).
\]
Its dimension has minimal polynomial
\[
 z^4-2304z^3+184320z^2-4718592z+37748736,
\]
whose constant term is $2^{22}3^2$.

The Ising character rows are $(1,1,\sqrt2)$, $(1,1,-\sqrt2)$ and
$(1,-1,0)$, with formal codegrees $4,4,2$. Thus for a character $\chi$
of $B$, let $f_B$ be the product of these codegrees and write
$a=\chi(A)=u+v\sqrt2$. If $\chi(L)=-1$, the identity $A=LA$ forces
$a=0$; its two extensions have formal codegree $2f_B$. If
$\chi(L)=1$, their formal codegrees are
\[
 f_\pm=\frac{f_B}{2}\left(4+a^2\pm a\sqrt{a^2+4}\right).
\]
Every $D/f_\pm$ is an algebraic integer and every $f_\pm$ belongs to
$K$ \cite[Corollaries 2.14--2.15]{Ostrik}.

First, $f_+f_-=f_B^2(a^2+4)$ implies that
\[
 n(u,v)=(u^2+2v^2+4)^2-8u^2v^2
\]
has no prime divisor other than $2$ and $3$. Second, the primitive
minimal polynomial of $D/f_\pm$ must be monic. Finally, if $a\ne0$,
field membership implies $\sqrt{a^2+4}\in K$. A square root in
$K=\mathbb Q(\sqrt2)\oplus\sqrt3\mathbb Q(\sqrt2)$ of an element of
$\mathbb Q(\sqrt2)$ lies in one of those two summands. Taking its
quadratic norm shows that $n(u,v)$ must be a square in $\mathbb Q$,
hence an integer square.

The finite domain is particularly small. If $z$ is the coefficient of
$\sigma\otimes\sigma$ in $A$, then $0\le z\le3$. The coefficients
of $1\otimes1,g\otimes g$ sum to $3-z$, those of
$1\otimes g,g\otimes1$ also sum to $3-z$, and the coefficients of the
four $\sigma$-supported simples sum to four. This gives $1050$ possible
$A$ for $L=1$ and six satisfying $A=eA$ for $L=e$.
The exact checker produces
\[
1056\xrightarrow{\text{prime norms}}158
\xrightarrow{D/f\text{ integral}}130
\xrightarrow{\text{dimension field}}58.
\]
Exchanging the two Ising factors, independently multiplying either
$\sigma$ generator by the fermion in the other factor, and replacing
$t$ by $vt$ for invertible $v\in B$ partition these alternatives into
fourteen orbits. These are sufficient representatives; no assertion
that distinct listed orbits are categorically distinct is needed.
\end{proof}

The script \texttt{verify\_relative\_reduction.py} checks the original
tensor, reconstructs all domains, and writes the exact witnesses and
representatives to \texttt{relative\_reduction\_certificate.json}.
It runs under \texttt{python -O}, without numerical optimizers.
The relative center need not be braided, its simples need not lift to
the full center, and the surviving rings are not asserted to have
categorifications. This replaces the present structural question by a
small family of rank-eighteen fusion-ring targets.

\begin{thebibliography}{9}
\bibitem{DGNO} V.~Drinfeld, S.~Gelaki, D.~Nikshych, V.~Ostrik,
\emph{On braided fusion categories I}, Selecta Math. 16 (2010), 1--119,
\href{https://arxiv.org/abs/0906.0620}{arXiv:0906.0620}.
\bibitem{GNN} S.~Gelaki, D.~Naidu, D.~Nikshych,
\emph{Centers of graded fusion categories}, Algebra Number Theory 3
(2009), 959--990, \href{https://arxiv.org/abs/0905.3117}{arXiv:0905.3117}.
\bibitem{Ostrik} V.~Ostrik, \emph{Pivotal fusion categories of rank 3},
Mosc. Math. J. 15 (2015), 373--396,
\href{https://arxiv.org/abs/1309.4822}{arXiv:1309.4822}.
\end{thebibliography}
\end{document}
